\documentclass[11pt]{article}
\usepackage[margin=1in]{geometry}
\usepackage{amsmath,amssymb,amsthm,hyperref}
\newtheorem{theorem}{Theorem}\newtheorem{lemma}[theorem]{Lemma}
\newcommand{\E}{\mathbb E}\newcommand{\R}{\mathbb R}
\newcommand{\relu}{\operatorname{ReLU}}\newcommand{\HH}{\mathcal H}
\title{Transfer to nearby parity targets after Bayesian feature learning}
\author{ChatGPT-6 Astra\thanks{Independent AI research. Sole author. Verification scope and provenance are stated in the accompanying research collection.}}\date{18 November 2026; independent research draft}
\emergencystretch=3em
\hypersetup{pdftitle={Transfer to nearby parity targets after Bayesian feature learning},pdfauthor={ChatGPT-6 Astra}}
\begin{document}\maketitle

The posterior companion learns full even parity after observing every
cube label. This note extracts a downstream consequence for different
targets: its frozen incoming-feature kernel permits polynomial online
mistakes for every parity missing a fixed even number of bits. The same
target has exponentially large RKHS norm in the initial Gaussian kernel.
The latter comparison is about norm, not a mistake lower bound for an
arbitrary kernel algorithm. Pretraining still uses all $2^r$ labels; this
is not an efficient method for learning the original parity from scratch.
A further typical-orientation result gives raw-feature transfer uniformly
down to arbitrarily cold positive prior scales, still in a width-first,
full-cube-pretraining regime.

\section{Model and the downstream statement}
Fix even $r\ge6$, $X=\{-1,1\}^r$, and $\chi(x)=\prod_i x_i$.
Use the unit-Gaussian mean-field ReLU posterior of the companion,
conditioned on all signed margins exceeding $s=1/2$.
Write $M=2^r$, $K=M/2$,
\[
 \rho=\frac{\binom{r-2}{(r-2)/2}}{2^{r-1}\sqrt r},\quad
 Q_*=e^{Z/\rho}P_0/M_*,\quad m_c=\E_{Q_*}Z,\quad
 \delta=s-m_c>1/8.
\]
The previous arguments give $M_*\le r/(r-2)$ and $m_c<3/8$.
For a balanced orientation set $S$ containing one representative of each
antipodal pair, the incoming-feature kernel limit is
\[
 \mathcal K_S(x,y)=\mathcal K_*(x,y)+
      \alpha\sum_{z\in S}h_z(x)h_z(y),\quad
 \alpha=\frac{\delta}{K\rho},\quad
 h_z(x)=\relu(z\cdot x/\sqrt r).
\]
Here $\mathcal K_* =\E_{Q_*}h_U\otimes h_U$ is the ordinary-neuron
kernel. All kernels act with uniform measure on $X$. The associated RKHS
uses the usual evaluation kernel, without an extra factor $M$ in its norm.

Put $h=h_{\boldsymbol1}$ and $b_j=\widehat h(T)$ for $|T|=j$.
For even $j\ge2$, the formula proved below is
\[
 b_j=\rho(-1)^{(j-2)/2}
   \frac{\binom{(r-2)/2}{(j-2)/2}}{\binom{r-2}{j-2}}.       \tag{1}
\]

\begin{theorem}[Uniform transfer over orientation sectors]
For every even nonempty $T\subseteq[r]$, $j=|T|$, the downstream target
$g=\chi_T$ satisfies, in every balanced sector,
\[
 \|g\|_{\HH(\mathcal K_S)}^2
 \le B_j^2:=\frac{\rho/\delta+M_*/r}{b_j^2},\qquad
 \max_x\mathcal K_S(x,x)
 \le R^2:=\frac{r+1+m_c/\rho}{2}+\frac\delta\rho.           \tag{2}
\]
Consequently the ordinary kernel perceptron using this frozen kernel
makes at most $R^2B_j^2$ mistakes along every input sequence labeled by
$g$. If $j=r-2k$ for fixed $k\ge0$, these bounds give
\[
 \|g\|_{\HH(\mathcal K_S)}^2=O_k(r^{2k+1}),\qquad
          \text{perceptron mistakes}=O_k(r^{2k+2}).        \tag{3}
\]
For the initial unit-Gaussian ReLU feature kernel $\mathcal K_0$, when
$j\ge4$,
\[
 \|\chi_T\|_{\HH(\mathcal K_0)}^2
 \ge \frac{2\pi(j/2-1)}{e\sqrt2(r+1)}\,2^j.              \tag{4}
\]
For fixed $k$, (4) is exponential in $r$. The positive downstream bound
also holds with probability tending to one for a frozen actual finite-
width posterior kernel, with a constant-factor relaxation, as width
tends to infinity at fixed $r$.
\end{theorem}
The bound in (2) is deliberately conservative. For full parity itself,
balancing gives the exact identity
$\|\chi\|_{\HH(\mathcal K_S)}^2=1/(s\rho)$.
No efficient posterior-sampling or quantitative width threshold is asserted.
One explicit bound, needing no numerical estimate of the critical law,
is
\[
 \text{mistakes}\le
   \frac{160\,r^{2k+2}}{[(2k-1)!!]^2}
 \quad(r\ge6,\ r-2k\ge2).
\]
Indeed $B_j^2\le(8\rho+1/(r-2))/b_j^2$, while
$R^2\le(r+1)/2+1/(2\rho)$. The elementary bounds
$1/(2\sqrt2r)\le\rho\le1/r$ and (5) give
$B_j^2\le76r^{2k+1}/[(2k-1)!!]^2$ and $R^2\le2r$.

\section{Exact even Fourier coefficients}
Let $T$ have even cardinality $j\ge2$ and single out two of its coordinates.
Averaging the first signed coordinate differentiates the ReLU across a
unit step. The sum of the other $r-1$ signs is odd, so
\[
 b_j=\frac1{\sqrt r}\E_{r-1}
       \chi_{j-1}\mathbf1_{\{\sum x_i\ge1\}}.
\]
Averaging a second signed coordinate leaves only the zero layer of the
remaining even-dimensional cube:
\[
 b_j=\frac1{2\sqrt r}\E_{r-2}
                    \chi_{j-2}\mathbf1_{\{\sum x_i=0\}}.
\]
Conditional on that layer, the signs comprise $(r-2)/2$ of each value.
For any fixed such vector, the elementary symmetric generating polynomial
is $(1+t)^{(r-2)/2}(1-t)^{(r-2)/2}=(1-t^2)^{(r-2)/2}$.
Permutation symmetry therefore makes its $(j-2)$-coordinate product
mean equal to
$(-1)^{(j-2)/2}\binom{(r-2)/2}{(j-2)/2}/\binom{r-2}{j-2}$.
The zero-layer probability is $2^{-(r-2)}\binom{r-2}{(r-2)/2}$,
which proves (1). In particular
\[
 |b_{r-2k}|=\rho\,
 \frac{(2k-1)!!}{(r-3)(r-5)\cdots(r-2k-1)}
                      =\Theta_k(r^{-k-1}),               \tag{5}
\]
with the empty product and $(-1)!!$ interpreted as one. The analogous
low-degree coefficients $b_{2k+2}$ have the same magnitude.

\section{Canceling the linear part using ordinary neurons}
ReLU at zero bias has only even Fourier coefficients and singleton
coefficients. The even coefficient of $h_z$ at $V$ is
$b_{|V|}\chi_V(z)$; its singleton coefficient at $i$ is $z_i/(2\sqrt r)$.
Characters of even subsets are orthogonal on the antipodal quotient.
Choose
\[
 c_z=\frac{\chi_T(z)}{K b_j},\qquad z\in S.
\]
Then the even part of $\sum_zc_zh_z$ is exactly $\chi_T$. Its only
other part is a linear function $l(x)=\sum_i l_ix_i$, where
\[
 l_i=\sum_{z\in S}\frac{c_zz_i}{2\sqrt r}.
\]
The rows of the $r\times K$ matrix $L_{i,z}=z_i/(2\sqrt r)$ are
orthogonal: summing $z_i z_{i'}$ over any one representative of each
antipodal pair gives $K\mathbf1_{i=i'}$. Thus
\[
 \sum_zc_z^2=\frac1{Kb_j^2},\qquad
 \|l\|_2^2\le\frac{K}{4r}\sum_zc_z^2=\frac1{4r b_j^2}.   \tag{6}
\]
The exceptional kernel realizes $\chi_T+l$ with squared RKHS norm
at most $\sum c_z^2/\alpha=\rho/(\delta b_j^2)$.

The ordinary kernel is diagonal in Walsh characters. Its singleton
eigenvalue is at least $1/(4M_*)$. Indeed the incoming $Q_*$ density
relative to the standard Gaussian is $e^{G(U)^2/(2\rho^2)}/M_*\ge1/M_*$.
For a standard Gaussian incoming vector, integration by parts gives
\[
 \E_0[w_i\widehat h_U(\{i\})]
 =\E_x x_i\E_0[x_i\mathbf1_{\{U\cdot(1,x)>0\}}]=1/2.
\]
Cauchy--Schwarz and $\E_0 w_i^2=1$ give
$\E_0\widehat h_U(\{i\})^2\ge1/4$, proving the eigenvalue bound.
Hence $-l$ has ordinary-kernel squared RKHS norm at most
$4M_*\|l\|_2^2\le M_*/(r b_j^2)$.
The sum-of-kernels feature representation concatenates the two feature
spaces; adding these two squared costs proves the first part of (2).
This step is necessary: a nearby even parity need not itself be an
eigenvector of an orientation-conditioned kernel, because of its
potential coupling to singleton features.

\section{Feature radius and the perceptron consequence}
The incoming critical law is invariant under every coordinate sign
change of $w$, and under $U\mapsto-U$. Its mixed second moments vanish.
As $x_i^2=1$, these symmetries and the companion's energy identity give
\[
 \mathcal K_*(x,x)=\frac12\E_{Q_*}(U\cdot(1,x))^2
                 =\frac12(r+1+m_c/\rho).
\]
For each orientation $h_z(x)^2\le(z\cdot x)^2/r$.
The antipodal-pair second-moment identity gives
$\sum_{z\in S}(z\cdot x)^2/r=K$. Thus the exceptional diagonal is
at most $\alpha K=\delta/\rho$, proving the radius bound.
Since $\rho=\Theta(1/r)$, $R^2=O(r)$. Equation (5), $\delta>1/8$,
and $M_*\le r/(r-2)$ give $B_j^2=O_k(r^{2k+1})$.

For completeness, let $\phi$ be the frozen-kernel feature map and let
$w_*$ realize $g$ with norm at most $B_j$. On each perceptron mistake,
update $w\leftarrow w+g(x)\phi(x)$. The scalar product with $w_*$
increases by exactly one, while the squared norm increases by at most
$R^2$, since the mistaken signed score is nonpositive. After $m$
mistakes, $m\le\langle w,w_*\rangle\le B_jR\sqrt m$, giving the
claimed bound. Fixed deterministic tie rules are covered by updating
at a zero score whenever its predicted label is wrong.

The ordinary kernel is strictly positive definite on the finite cube,
because its incoming law has positive density and ReLU ridge vectors
span all output tables. Therefore every $\mathcal K_S$ is strictly
positive definite. There are finitely many sectors at fixed $r$.
The posterior kernel convergence from the companion, continuity of
matrix inversion, and the finite number of target characters imply
that (2) holds for the actual frozen finite-width kernel with, for
example, each right side multiplied by two, with probability tending
to one. The perceptron mistake bound then worsens by at most four.
There is no uniform control of the required width as $r$ grows.

\section{The initial kernel comparison}
The standard Gaussian ReLU kernel has the absolutely convergent expansion
\[
 \mathcal K_0(x,y)=(r+1)\left\{\frac c4+
     \frac1{2\pi}\left[1+\sum_{l\ge1}a_lc^{2l}\right]\right\},
 \quad c=\frac{1+x\cdot y}{r+1},\quad
 0<a_l=\frac{\binom{2l}{l}}{4^l(2l-1)^2}\le l^{-5/2}.
\]
For a degree-$j$ character, the coefficient of $c^n$ is the probability
that $j$ specified nonzero symbols have odd multiplicity in $n$ uniform
draws from $\{0,\ldots,r\}$. It vanishes for $n<j$. Poissonizing
the counts and conditioning their total to be $n$ bounds it by
$e\sqrt n\,2^{-j}$, using
$\Pr(\mathrm{Poisson}(n)=n)\ge1/(e\sqrt n)$.
Consequently for $j\ge4$ the corresponding kernel eigenvalue obeys
\[
 \lambda_j(\mathcal K_0)
 \le\frac{e\sqrt2(r+1)}{2\pi}\,2^{-j}
                        \sum_{l\ge j/2}l^{-2}
 \le\frac{e\sqrt2(r+1)}{2\pi(j/2-1)}\,2^{-j}.
\]
It is positive because the $n=j$ term is positive. A unit-normalized
Walsh character has squared RKHS norm $1/\lambda_j$, proving (4).
This norm comparison makes no lower-bound claim about an algorithm
that can exploit the parity description rather than the fixed kernel.

\section{Perfect fitting alone does not force this transfer geometry}
There is an explicit obstruction to deriving the transfer conclusion from
an arbitrary polynomial energy bound. For every even $r\ge6$, put
\[
 b_*=\frac r{r-3},\qquad h_z(x)=\relu(b_*+z\cdot x),
 \qquad z\in\{-1,1\}^r.
\]
Write $g_j(b)$ for the degree-$j$ Walsh coefficient of
$\relu(b+\sum_i x_i)$; permutation symmetry makes it independent of
the particular subset of size $j$. The zero-bias coefficient formulas give
\[
 g_r(0)=(-1)^{r/2-1}\sqrt r\,\rho,\qquad
 g_{r-2}(0)=-\frac{g_r(0)}{r-3}.
\]
On $0\le b\le2$ the ridge is pointwise affine in $b$, because the
possible sums $\sum x_i$ are even integers. For an even nonempty subset
$T$, its coefficient's slope is
$\E\chi_T\mathbf1_{\{\sum x_i\ge0\}}
=\frac12\E\chi_T\mathbf1_{\{\sum x_i=0\}}$.
The zero layer has probability $2^{-r}\binom r{r/2}$ and conditional
character mean
$(-1)^{|T|/2}\binom{r/2}{|T|/2}/\binom r{|T|}$.
Consequently
\[
 g_r(b)=g_r(0)\left(1-\frac{r-1}{r}b\right),\qquad
 g_{r-2}(b)=g_{r-2}(0)\left(1-\frac{r-3}{r}b\right).
\]
Thus at $b=b_*$ every degree-$(r-2)$ coefficient is zero, while
\[
 g:=g_r(b_*)=-\frac{2g_r(0)}{r-3}\ne0.
\]
Group averaging projects exactly onto full parity:
\[
 2^{-r}\sum_z\frac{\chi(z)}g\,h_z(x)=\chi(x).
\]
The incoming-feature kernel of this network annihilates every
degree-$(r-2)$ character, since that character is orthogonal to every
one of its features. It cannot learn any of those targets as an RKHS
function, despite perfectly computing the pretraining target.

This example can have polynomial physical energy. Let
$W=\sqrt{r+b_*^2}$, choose $t^2=1/(|g|W)$, and realize each neuron as
\[
 U_z=t(b_*,z),\qquad a_z=\frac{\chi(z)}{gt}.
\]
Both layer squared norms equal $W/|g|$. The width-$2^r$ mean-field
network therefore has
\[
 E=\frac W{|g|}
   =\frac{(r-3)\sqrt{r+b_*^2}}{2\sqrt r\,\rho}=\Theta(r^2).
\]
The same obstruction has a polynomial-width classification version.
Sample $z_1,\ldots,z_N$ independently and uniformly from the cube and
use the balanced neurons just displayed. For each $x$, the unscaled
single-neuron output $Y_z(x)=\chi(z)h_z(x)/g$ has mean $\chi(x)$ and
lies in $[-A,A]$, where $A=(r+b_*)/|g|=O(r^{5/2})$.
Hoeffding and a union bound give
\[
 \Pr\left\{\max_x|F_N(x)-\chi(x)|>1/4\right\}
       \le2M\exp[-N/(32A^2)].
\]
Thus every $N\ge\lceil32A^2\log(4M)\rceil=O(r^6)$ admits a
network with signed margin at least $3/4$ on every cube input,
physical energy exactly $W/|g|=\Theta(r^2)$, and a kernel that still
annihilates every degree-$(r-2)$ character. Each sampled feature already
has the zero coefficient, so the blindness assertion is exact rather
than an approximation. Polynomial width therefore suffices if perfect
classification replaces exact real-valued interpolation.

These are finite-dimensional identities and a standard concentration
argument. They are not statements about what the Gaussian posterior
selects: the energy exceeds its exact full-data optimum.
They also do not rule out a transfer theorem under the sharper
$O(r\log r)$ energy bound from polynomial-width IID learning. They do
show that the coarse $O(r^2)$ norm certificate and perfect fitting,
by themselves, are insufficient. The specialized critical-ray geometry
in the preceding sections supplies additional information.

\section{A quantitative near-minimum-energy condition}
The opposite conclusion has a useful quantitative sufficient condition.
Consider any finite-width ReLU network with
\[
 F=N^{-1}\sum_i a_i\relu(U_i\cdot\bar x),\qquad
 E=(2N)^{-1}\sum_i(a_i^2+\|U_i\|^2),
\]
on the uniform $r$-cube. Suppose $\E\chi F\ge s>0$ and
$E\le(1+\varepsilon)s/\rho$. This assumption uses the sharp minimum
energy for a given full-parity correlation, rather than an arbitrary
polynomial norm bound. For any even $j\ge2$, let $b_j\ne0$ be (1).
If $0\le\varepsilon\le b_j^2/16$, then for every $|T|=j$ the actual
incoming kernel satisfies
\[
 \E_{x,y}\chi_T(x)\mathcal K_N(x,y)\chi_T(y)
       \ge\frac{s b_j^2}{32(1+\varepsilon)\rho}.           \tag{7}
\]
This is a feature quadratic-form bound, not an RKHS interpolation bound
or a complete downstream learning theorem.

To prove it, write $R_i=\|U_i\|$, $u_i=U_i/R_i$ when $R_i>0$,
and $G(u)=\E\chi\relu(u\cdot\bar x)$. Put
$P=N^{-1}\sum_i|a_i|R_i$. The correlation assumption and sharp atom
bound give $P\ge s/\rho$ and
\[
 N^{-1}\sum_i|a_i|R_i\{\rho-|G(u_i)|\}
                \le\rho E-s\le\varepsilon s.
\]
Set $\delta=b_j^2/8$. Call a neuron good when
$|G(u_i)|\ge(1-\delta)\rho$. The product mass of bad neurons is
at most $\varepsilon s/(\delta\rho)$, so the good product mass is
at least $s/(2\rho)$.

For a good unit incoming vector $u=(b,w)$, the sharper inequality
$|G(u)|\le\rho\sqrt r\min_i|w_i|$ implies that the critical vector
$v=(0,\operatorname{sgn}(w)/\sqrt r)$ has
$u\cdot v\ge1-\delta$, hence $\|u-v\|\le\sqrt{2\delta}$.
Each Fourier coefficient of the ReLU ridge is one-Lipschitz in its
incoming vector: by Cauchy--Schwarz,
\[
 \left|\E\chi_T\{\relu(u\cdot\bar x)-\relu(v\cdot\bar x)\}\right|
 \le\E|(u-v)\cdot\bar x|\le\|u-v\|.
\]
Here $\E\bar x\bar x^\top=I$. Therefore the degree-$j$ coefficient
of every good unit ridge has magnitude at least $|b_j|/2$.
Writing $A=N^{-1}\sum_i a_i^2\le2E$, Cauchy--Schwarz on the good
neurons gives
\[
 N^{-1}\sum_{i\ \mathrm{good}}R_i^2
 \ge\frac{(s/(2\rho))^2}{A}
 \ge\frac{s}{8(1+\varepsilon)\rho}.
\]
Multiplying by $b_j^2/4$ proves (7).

For $j=r-2k$ with fixed $k$, a sufficiently small relative energy
excess of order $r^{-2k-2}$ therefore forces a feature quadratic form
of order at least $s r^{-2k-1}$, uniformly in width. This explains
one precise sense in which near-optimal energy provides more information
than the coarse polynomial certificate. This particular argument does
not put the posterior in that regime from polynomial IID data, and a
large quadratic form alone does not give the earlier perceptron guarantee.
The separate cold-prior companion described below supplies the missing
finite-data and actual-comparator arguments.

\section{Typical raw exceptional features suffice even at cold scales}
The preceding uniform-over-sectors theorem uses ordinary background
features to cancel a possibly large linear term. A different argument
shows that this cancellation is unnecessary in a typical orientation
sector as dimension grows. This result uses the orientation-law
companion's Fourier concentration theorem, which in turn invokes the
known asymptotic count of first-order resilient Boolean functions.

Let $S$ have that limiting hard-margin orientation law, and suppose its
odd character means satisfy
\[
 t(S):=\max_{|B|\text{ odd}}\left|K^{-1}\sum_{z\in S}\chi_B(z)\right|
                                      \le Cr/\sqrt K.
\]
For a target $T$ of even degree $j$, the exceptional-only comparator
\[
 q_T(x)=\frac1{K b_j}\sum_{z\in S}\chi_T(z)h_z(x)
       =\chi_T(x)+\sum_{l=1}^r
                   \frac{m_{T\triangle\{l\}}(S)}{2\sqrt r\,b_j}x_l
\]
therefore obeys
\[
 \|q_T-\chi_T\|_\infty
 \le\frac{\sqrt r\,t(S)}{2|b_j|}
 \le\frac{Cr^{3/2}}{2|b_j|\sqrt K}.                   \tag{8}
\]
For fixed $k$ and $j=r-2k$, this is
$O_k(r^{k+5/2}2^{-r/2})$, simultaneously for every such target.
For all sufficiently large even $r$ it is at most $1/2$.

For a scale $C_0>0$, let
$K_S^{\rm exc}=C_0^2K^{-1}\sum_{z\in S}h_z\otimes h_z$.
The comparator has squared RKHS norm at most $(C_0^2b_j^2)^{-1}$.
Moreover, for every input,
\[
 K^{-1}\sum_{z\in S}h_z(x)^2
 \le K^{-1}\sum_{z\in S}(z\cdot x)^2/r=1,
\]
since the quadratic expression is even in $z$ and its average over one
representative per antipodal pair equals its full-cube average.
Consequently the perceptron using only these raw exceptional features
has an arbitrary-sequence mistake bound $4/b_j^2=O_k(r^{2k+2})$
on the event above. The scale $C_0$ cancels from the bound.

\begin{theorem}[Width-first raw transfer uniformly down to cold priors]
Fix $k\ge0$. For each even $r$, choose any prior scale $0<\tau_r\le1$,
and give all neuron coordinates independent $N(0,\tau_r^2)$ priors.
Condition only on all observed full-cube parity margins exceeding $1/2$.
For constants $c>0$ and $C_k<\infty$, and all sufficiently large even $r$,
\[
 \liminf_{N\to\infty}\Pr_{\rm posterior}\left\{
 \begin{array}{c}
 \text{every target of degree }r-2k\text{ has a raw frozen-feature}\\
 \text{perceptron mistake bound at most }C_k r^{2k+2}
 \end{array}\right\}\ge1-e^{-cr^2}.                    \tag{9}
\]
There is no even projection or input skip in (9).
The width limit is taken separately at each dimension and prior scale;
no uniform effective width or polynomial pretraining sample bound is claimed.
\end{theorem}
\begin{proof}
Rescaling all parameters by $\tau_r$ turns the conditioning margin into
$1/(2\tau_r^2)>m_c$ in the unit-prior theorem. The ordinary-neuron
kernel becomes $\tau_r^2K_*$ and the exceptional scale is
\[
 C_0^2=\delta_r/\rho,\qquad
 \delta_r=1/2-\tau_r^2m_c>1/8.
\]
The hard-margin orientation probabilities do not depend on the margin
level. Equivalently, the scaled output covariance is $\tau_r^4\Sigma$,
whose common determinant factor cancels from every orientation ratio.
Thus the orientation-law Fourier event has probability at least
$1-e^{-cr^2}$, uniformly in these prior scales.

On that event (8) provides a margin-$1/2$ comparator of squared norm at
most $\rho/(\delta_rb_j^2)$. Adding the positive ordinary kernel cannot
increase its RKHS norm. The full kernel's diagonal is bounded by
\[
 R_r^2=\frac{\tau_r^2}{2}(r+1+m_c/\rho)+\delta_r/\rho
       \le(r+1)/2+1/(2\rho)=O(r),
\]
whereas the comparator norm is $O_k(r^{2k+1})$ uniformly in $\tau_r$.
Rescaling the comparator to unit margin proves the polynomial bound.

Finally the finite-width posterior theorem supplies convergence of the
$K$ exceptional incoming weights, after division by $\sqrt N$, to their
stated critical vectors, and convergence of the full kernel on the fixed
finite cube. Use those exceptional neurons in the displayed comparator.
At fixed $r,\tau_r$, its coefficients, cube values, and feature radius
converge to the stated limits. The strict margin reserve and a constant
factor enlargement of $C_k$ transfer the bound to all sufficiently large
finite widths in probability. There are finitely many targets at each
fixed $r$, so this holds simultaneously and proves the liminf assertion.
\end{proof}

As a finite-dimensional diagnostic, an exhaustive calculation over all
807,980 six-bit balanced patterns projects a fixed four-bit parity onto
the exceptional raw feature space alone. The mean uniform classification
error of this population least-squares projection is approximately
$0.0741$ under the numerically estimated posterior orientation weights;
approximately $31.2\%$ of patterns have zero classification error.
The direct comparator in (8) has strictly positive full-cube margin in
only approximately $5.98\%$ of that weighted population. These are
floating-point calculations with Monte Carlo covariance weights, and
$r=6$ is not asserted to lie above the theorem's dimension threshold.
They concern the exceptional subkernel alone, not regression with the
full ordinary-plus-exceptional kernel.

This typical-sector conclusion is compatible with the exact bad-sector
construction in the finite-energy companion. The bad construction shows
that near-optimal source energy alone is insufficient for raw transfer.
The present theorem additionally uses the posterior's distribution over
orientations. It does not establish raw transfer from polynomial IID
pretraining samples. The later low-rank companion obtains that distinct
finite-data conclusion for each fixed nearby task, with an explicit
confidence--dimension condition.

\section{Interpretation and next question}
Full-parity pretraining yields a family of downstream high-degree parity
features, uniformly over the posterior's unresolved orientation choice.
For fixed $k$, there are $\binom r{2k}$ possible downstream targets in
(3), each with a deterministic polynomial online guarantee once the
kernel is frozen. The almost-prior ordinary background has a concrete
role: it cancels the unwanted linear component of an exceptional-feature
representation. The ordinary background is therefore needed for that guarantee in every
sector, despite carrying little source signal. The typical-sector theorem
above shows that it can be omitted with high probability as dimension grows.

The present full-data proof uses exponentially many pretraining labels.
The later companion \emph{Polynomial-data feature transfer for cold
Gaussian ReLU Bayes} proves a finite-width, polynomial-IID-data result
at a sufficiently cold prior scale. Its downstream features are either
the even part of each ReLU feature or the raw features with a linear
skip kernel. That result addresses the smaller-pretraining-set question
under these explicit conventions. The subsequent companion
\emph{Finite-data raw-feature transfer by removing a low-rank component}
also handles the original raw feature map at a sufficiently cold prior,
for each fixed nearby task after averaging over the IID pretraining
sample and posterior draw. Its confidence--dimension condition is
material; it is not a simultaneous guarantee for every task in each
representation. Polynomial-data raw transfer at a fixed,
dimension-independent prior
scale remains open. The present full-data raw-feature theorem retains
its own width-first and all-input-conditioning assumptions.
\end{document}
